\chapter{Conclusions}

\section{Summary}
This thesis studied how to detect very small human targets in aerial disaster
imagery with a model small enough to deploy. Starting from a YOLO11m
baseline, an additive ablation measured the effect of CBAM attention, a
high-resolution P2 detection scale, and a state-space neck on the C2A dataset
under one training and evaluation protocol. The P2 scale was the component
that mattered: it raised very-tiny recall from $0.743$ to $0.757$ and
$\mathrm{AP}_{50}$ from $0.843$ to $0.853$ at a cost of half a million
parameters and one millisecond of latency. CBAM was near-free, removing about
one million parameters while slightly improving the operational scores. The
state-space neck added $2.4$~million parameters and $2.8$ times the latency
without improving any accuracy metric, and was excluded on that evidence. The
recommended CBAM~+~P2 detector reaches $\mathrm{AP}_{50}$ $0.853$ and $F_2$
$0.844$ at $19.6$~million parameters and $14.6$~ms per image, second among
all published detectors on this benchmark and a third of the leader's size;
test-time augmentation at $1280$~px lifts very-tiny recall further, from
$0.758$ to $0.850$, at $60$~ms per image without retraining. The whole
training campaign was carbon-accounted at $9.0$~kg of CO$_2$.

An extended phase then hardened and stretched these results. A split audit
found 98\% scene sharing in the official benchmark and re-based everything on
a scene-disjoint split, at a measured cost of $1.6$ points of
$\mathrm{AP}_{50}$ and no changed conclusion. With the architecture search
stalled, the productive lever moved to training supervision: a Wasserstein
blend inside the label assigner raised sub-8-pixel recall by $2.08$ points
with a bootstrap interval that excludes zero, while a frequency-gated
evidence branch failed its paired test and joined the state-space neck as a
documented negative. Real footage collected for this thesis completed the
picture. The synthetic-trained model transfers poorly zero-shot
($\mathrm{AP}_{50}$ $0.335$ on drone footage, $0.133$ on real disaster
frames), a few hundred labelled frames close most of the drone gap ($0.795$)
and triple the within-video disaster score ($0.411$) at almost no cost to
the source benchmark, and a disaster event held out from all training shows
where that adaptation stops: there the gain is not statistically
significant. Twice over, supervision moved the metric where added
architecture did not.

\section{Limitations}
The study has clear limits, summarised here so the results are read at their
proper weight.
\begin{itemize}
  \item \textbf{Semi-synthetic primary benchmark.} The main results are on
        C2A, whose scenes composite real backgrounds with segmented human
        figures, and whose paste-based annotations cap localisation accuracy
        for every model compared. The extended phase measured the transfer
        gap directly instead of leaving it open, and the gap is large:
        zero-shot $\mathrm{AP}_{50}$ falls from $0.844$ on C2A to $0.335$ on
        the drone test and $0.133$ on real disaster frames.
  \item \textbf{Small real evaluation sets.} The real test sets hold 23 to
        60 frames each, so their scores carry wide intervals and should be
        read as evidence of direction, not precise levels. The within-video
        disaster result also shares source videos between training and
        evaluation frames, deliberately and with de-duplication; the honest
        cross-event number is the held-out one, and it is not significant.
  \item \textbf{Sub-8-pixel people in real footage remain undetected.} On
        the disaster R-set, recall below 8~px is zero for every model
        tested, across all 54 such instances. The regime the thesis targets
        on synthetic data is still open on real data, and the adaptation
        fine-tune also paid 13 points of precision on the held-out event.
  \item \textbf{Single seed per training run.} The ablation and the extended
        runs use seed 0. The bootstrap intervals of
        Section~\ref{sec:stats} treat test-set variation only; replicate
        runs of the assignment pair at two further seeds were in progress at
        the time of writing.
  \item \textbf{Sparse medium-target evidence.} The medium band holds only
        317 of roughly 72{,}500 test instances on the official split (410 on
        the scene-disjoint split), so its recall estimates are noisy.
  \item \textbf{No on-device measurement.} Latency was measured on a desktop
        GPU. Embedded or airborne hardware will behave differently, and no
        field trial with a drone was performed.
\end{itemize}

\section{Recommendations for Future Works}
The limitations above set the agenda, and the extension phase has already
converted two earlier recommendations into results: the enhanced dataset
with the author's own imagery now exists as the real-imagery suite of
Table~\ref{tab:owndata}, and its adaptation study is reported in
Section~\ref{sec:realresults}. What follows is what remains.
\begin{itemize}
  \item \textbf{Multi-event disaster data.} The cross-event result shows
        that 60 frames from two events adapt the model to those events, not
        to the disaster domain. A disaster training set spanning many
        events, rather than many frames per event, is the most direct route
        to generalisation, and the annotation pipeline built here makes each
        added event cheap.
  \item \textbf{Completing the seed replication.} The replicate runs of the
        assignment pair should be finished and folded into
        Table~\ref{tab:d2pair} as a mean and spread, closing the last gap in
        the statistical story.
  \item \textbf{Server-side deployment.} Work is under way to serve the
        recommended model behind a streaming pipeline in which a drone
        transmits footage to a ground server that runs detection and returns
        flagged frames to the operator. This keeps the airborne payload
        light, allows the $1280$-px test-time augmentation mode when the
        link budget permits, and is the practical route to field use.
  \item \textbf{Composing the levers.} The assignment blend, the P2 scale,
        and inference-time slicing all target the same sub-8-pixel band by
        different mechanisms; whether they stack on real footage, where that
        band is still lost entirely, is an open and testable question.
  \item \textbf{Validation on public rescue imagery.} Zero-shot and
        fine-tuned evaluation on collections such as
        SARD~\cite{sambolek2021sard}, HERIDAL~\cite{bozic2019heridal},
        VisDrone~\cite{zhu2021visdrone}, and
        Okutama-Action~\cite{barekatain2017okutama} would ground the
        generalisation claim beyond the author's own sets.
  \item \textbf{A purpose-built state-space design.} Two rejected branches
        do not close the architecture direction: a neck designed around the
        geometry of very small targets, rather than a generic scan or a
        generic frequency gate, remains worth investigating, now with a
        validated protocol and a statistical harness to test it against.
\end{itemize}
